\appendices

\section{Impact of DP-FedAvg on Accuracy, Convergence, and Fairness}
\label{app:dp_impact}

At federation round $k$, let
$\Delta_{j,k}=\phi_{R,k}^{j}-\phi_{R,k}^{\mathrm{global}}$ denote the
risk-model update of insurer $j$. The implemented clipped and perturbed
aggregation can be written as
\begin{align}
\bar{\Delta}_{j,k}
&=\Delta_{j,k}\min\!\left\{1,
\frac{C}{\|\Delta_{j,k}\|_2}\right\},\\
\phi_{R,k+1}^{\mathrm{global}}
&=\phi_{R,k}^{\mathrm{global}}
+\sum_{j\in\mathcal I}\alpha_{j,k}\bar{\Delta}_{j,k}
+\zeta_k,
\qquad
\zeta_k\sim\mathcal N(0,\sigma_{\mathrm{DP}}^2I),
\label{eq:dp_appendix_update}
\end{align}
where $\alpha_{j,k}$ is proportional to the local sample size. Since
$\mathbb E[\zeta_k]=0$ but
$\mathbb E\|\zeta_k\|_2^2=d_R\sigma_{\mathrm{DP}}^2$, the perturbation does
not directly shift the mean update, but increases its variance. Consequently,
larger $\sigma_{\mathrm{DP}}$ can increase the prediction-error floor and
slow convergence. Clipping can additionally introduce bias whenever
$\|\Delta_{j,k}\|_2>C$.

We quantify predictive fairness through
\begin{equation}
\Delta_{\mathrm{ins}}
=\max_{j\in\mathcal I}\mathrm{MSE}_{j}
-\min_{j\in\mathcal I}\mathrm{MSE}_{j},
\qquad
\Delta_{\mathrm{grp}}
=\max_q\mathrm{MSE}_{q}-\min_q\mathrm{MSE}_{q},
\label{eq:dp_fairness_metrics}
\end{equation}
where $q$ denotes a patient latent-risk quartile. Lower values indicate more
uniform prediction quality across insurers or patient groups.

A controlled risk-learning ablation was performed using $100$ patients,
$4$ heterogeneous insurers, $120$ monthly rounds, federation every $10$
rounds, clipping threshold $C=1$, and $50$ common random seeds. The patient
risk function and stochastic state evolution follow the main simulation,
while the risk learner is evaluated separately from the price and coverage
networks to isolate the effect of the DP perturbation. We test
$\sigma_{\mathrm{DP}}\in\{0,0.01,0.05,0.10\}$, where
$\sigma_{\mathrm{DP}}=0$ is clipped FedAvg without Gaussian noise.

\begin{figure}[t]
\centering
\includegraphics[width=\columnwidth]{Figures/dp_accuracy.pdf}
\caption{Global risk-prediction MSE under different DP noise scales. Curves
show the mean over $50$ seeds and shaded regions show $95\%$ confidence
intervals. Periodic changes coincide with federation rounds.}
\label{fig:dp_accuracy}
\end{figure}

\begin{figure}[t]
\centering
\includegraphics[width=\columnwidth]{Figures/dp_convergence.pdf}
\caption{Empirical convergence month versus $\sigma_{\mathrm{DP}}$. Error
bars show $95\%$ confidence intervals over $50$ common seeds.}
\label{fig:dp_convergence}
\end{figure}

\begin{figure}[t]
\centering
\includegraphics[width=\columnwidth]{Figures/dp_fairness.pdf}
\caption{Insurer-level and patient-risk-group prediction disparities under
different DP noise scales. Error bars show $95\%$ confidence intervals.}
\label{fig:dp_fairness}
\end{figure}

\begin{table*}[t]
\centering
\caption{Effect of DP-FedAvg Noise on Risk Learning and Predictive Fairness
(Averaged Over 50 Seeds)}
\label{tab:dp_impact}
\renewcommand{\arraystretch}{1.12}
\begin{tabular}{ccccc}
\toprule
$\sigma_{\mathrm{DP}}$ & Final risk MSE & Convergence month &
$\Delta_{\mathrm{ins}}$ & $\Delta_{\mathrm{grp}}$ \\
\midrule
0.00 & $0.0013 \pm 0.0005$ & $25.2 \pm 30.1$ &
$0.0000 \pm 0.0000$ & $0.0028 \pm 0.0016$ \\
0.01 & $0.0025 \pm 0.0020$ & $36.8 \pm 40.4$ &
$0.0012 \pm 0.0020$ & $0.0029 \pm 0.0017$ \\
0.05 & $0.0284 \pm 0.0496$ & $68.0 \pm 40.2$ &
$0.0243 \pm 0.0498$ & $0.0134 \pm 0.0114$ \\
0.10 & $0.0925 \pm 0.1669$ & $72.1 \pm 37.7$ &
$0.0814 \pm 0.1671$ & $0.0349 \pm 0.0382$ \\
\bottomrule
\end{tabular}
\end{table*}

As shown in Figs.~\ref{fig:dp_accuracy}--\ref{fig:dp_fairness} and
Table~\ref{tab:dp_impact}, increasing the noise scale produces a clear
privacy--utility trade-off. Relative to clipped FedAvg, the final MSE rises
from $0.0013$ to $0.0284$ at $\sigma_{\mathrm{DP}}=0.05$, while the mean
convergence month increases from $25.2$ to $68.0$. The insurer disparity and
risk-group disparity also increase to $0.0243$ and $0.0134$, respectively.
Paired Wilcoxon tests with Holm correction confirm significant differences
from the non-noisy baseline at $\sigma_{\mathrm{DP}}\geq0.05$ for all four
metrics ($p<0.001$). At $\sigma_{\mathrm{DP}}=0.01$, the increase in final
MSE and insurer disparity is significant, whereas the changes in convergence
and group disparity are not. Therefore, DP noise does not inherently improve
fairness; under heterogeneous local data, strong perturbation can amplify
prediction disparities. Any improvement in the market-level outcomes of
FL-IMDR must consequently be attributed to the closed-loop regulator rather
than to improved risk-prediction accuracy.

The experiment reports the implemented noise scale directly. Mapping
$\sigma_{\mathrm{DP}}$ to a cumulative $(\varepsilon,\delta)$ guarantee
requires a specified neighboring relation and a privacy accountant across
all federation rounds.
